\paragraph{\texorpdfstring{$2^L$}{2^L}-进全息地址、有限 Grothendieck 秩障碍与 Lee--Yang 二次分辨率的接口压缩}
在单向量 \(2^L\)-进全息像集的精确自相似、地址动力系统的维数刚性、可消去幺半群的 Grothendieck 完成判据以及 Lee--Yang 三次数域链的共同三次基域都已确定之后，还可把这三条接口进一步压缩为下列五条直接后果。

\begin{theorem}[全息地址像的严格不交自相似分解]\label{thm:xi-time-part62ca-holographic-address-strict-selfsimilarity}
沿用结论 \ref{con:xi-time-part62b-rankone-hologram-selfsimilar-measure-dimension} 与结论 \ref{con:xi-time-part62c-godel-tate-address-fullshift-zeta-entropy} 的记号。设 \(E\) 为有限事件字母表，\(\mathrm{code}:E\hookrightarrow\{0,\dots,M\}\subset\ZZ\) 为单射，\(B=2^L>M\)，并记
\[
D:=\mathrm{code}(E),
\qquad
X(\mathbf e):=\sum_{t\ge 1}\mathrm{code}(e_t)\,B^{t-1}v
\qquad
(\mathbf e=(e_t)_{t\ge 1}\in E^{\NN}),
\]
\[
\mathcal X_E:=X(E^{\NN})\subset T_v.
\]
则 \(\mathcal X_E\) 满足严格的不交自相似分解
\[
\mathcal X_E
=
\bigsqcup_{a\in D}\bigl(av+B\mathcal X_E\bigr).
\]
更精确地，每一块 \(av+B\mathcal X_E\) 都是 \(\mathcal X_E\) 与某个 \(BT_v\)-陪集的交，因而在 \(\mathcal X_E\) 的相对拓扑下是 clopen。
\end{theorem}
\begin{proof}
由结论 \ref{con:xi-time-part62b-rankone-hologram-selfsimilar-measure-dimension} 的第 \(1\) 条，对 \(D=\mathrm{code}(E)\) 已有
\[
\mathcal X_D=\bigsqcup_{d\in D}\bigl(dv+B\mathcal X_D\bigr),
\]
且每一块都可写成 \(\mathcal X_D\) 与某个 \(BT_v\)-陪集的交。由于 \(\mathrm{code}\) 为单射，\(|D|=|E|\)，并且按定义 \(\mathcal X_E=\mathcal X_D\)。代换记号即得所述分解。证毕。
\end{proof}

\begin{theorem}[全息地址像的 Hausdorff 与 Minkowski 维数完全一致]\label{thm:xi-time-part62ca-holographic-address-exact-dimension}
沿用定理 \ref{thm:xi-time-part62ca-holographic-address-strict-selfsimilarity} 的记号，并在 \(\mathcal X_E\) 上赋予 \(B\)-进超度量
\[
d_B(x,y):=
\begin{cases}
0,& x=y,\\[2pt]
B^{-\nu_B(x-y)},& x\neq y,
\end{cases}
\qquad
\nu_B(z):=\max\{n\ge 0:\ z\in B^nT_v\}.
\]
则
\[
\dim_{\mathrm H}(\mathcal X_E)
=
\underline{\dim}_{\mathrm M}(\mathcal X_E)
=
\overline{\dim}_{\mathrm M}(\mathcal X_E)
=
\frac{\log |E|}{\log B}.
\]
\end{theorem}
\begin{proof}
由定理 \ref{thm:xi-time-part62ca-holographic-address-strict-selfsimilarity}，长度 \(n\) 的前缀给出 \(|E|^n\) 个两两不交的 cylinder，而每个 cylinder 都恰是某个 \(B^nT_v\)-陪集与 \(\mathcal X_E\) 的交，因此其直径正好为 \(B^{-n}\)。故尺度 \(B^{-n}\) 下的最少覆盖数满足
\[
N(B^{-n})=|E|^n.
\]
于是
\[
\underline{\dim}_{\mathrm M}(\mathcal X_E)
=
\overline{\dim}_{\mathrm M}(\mathcal X_E)
=
\lim_{n\to\infty}\frac{\log |E|^n}{-\log B^{-n}}
=
\frac{\log |E|}{\log B}.
\]

另一方面，结论 \ref{con:xi-time-part62c-godel-tate-address-ahlfors-regularity} 已给出 \(\mathcal X_E\) 上的等权 Bernoulli 推前测度是
\[
s:=\frac{\log |E|}{\log B}
\]
维 Ahlfors 正则测度，故由质量分布原理
\[
\dim_{\mathrm H}(\mathcal X_E)=s.
\]
与上式合并即得结论。证毕。
\end{proof}

\begin{corollary}[任意有限值连续观测都因子化于有限前缀]\label{cor:xi-time-part62ca-finite-valued-continuous-observables-prefix-factorization}
\leanverified{paper\_xi\_time\_part62ca\_finite\_valued\_continuous\_observables\_prefix\_factorization}
沿用定理 \ref{thm:xi-time-part62ca-holographic-address-strict-selfsimilarity} 的记号。设 \(A\) 为有限离散集，且
\[
F:\mathcal X_E\to A
\]
连续。则存在某个整数 \(n\ge 1\) 与映射
\[
\overline F:E^n\to A
\]
使得对任意 \(\mathbf e=(e_t)_{t\ge 1}\in E^{\NN}\) 都有
\[
F\!\bigl(X(\mathbf e)\bigr)=\overline F(e_1,\dots,e_n).
\]
等价地，存在映射
\[
\widetilde F:T_v/B^nT_v\to A
\]
使
\[
F=\widetilde F\circ \pi_n|_{\mathcal X_E},
\]
其中 \(\pi_n:T_v\to T_v/B^nT_v\) 为自然商映射。
\end{corollary}
\begin{proof}
由结论 \ref{con:xi-time-part62c-godel-tate-address-fullshift-zeta-entropy}，地址映射
\[
X:E^{\NN}\longrightarrow \mathcal X_E
\]
是同胚。故
\[
G:=F\circ X:E^{\NN}\to A
\]
连续。由于 \(A\) 离散，\(G^{-1}(a)\) 对每个 \(a\in A\) 都是 clopen。再因 \(E^{\NN}\) 的拓扑基由有限前缀 cylinder 构成，且 \(E^{\NN}\) 紧致，存在某个统一层数 \(n\) 使得每个长度 \(n\) 的 cylinder 都完全包含在某一条纤维 \(G^{-1}(a)\) 中。于是 \(G\) 在每个长度 \(n\) cylinder 上常值，可定义
\[
\overline F(e_1,\dots,e_n):=G(\mathbf e)
\]
其中 \(\mathbf e\) 为任意以 \((e_1,\dots,e_n)\) 为前缀的无限延拓；该定义与延拓无关。由 \(G=F\circ X\) 即得
\[
F\!\bigl(X(\mathbf e)\bigr)=\overline F(e_1,\dots,e_n).
\]

再由结论 \ref{con:xi-time-part62c-godel-tate-prefix-oracle-budget} 的第 \(1\) 条，长度 \(n\) 的前缀与 \(T_v/B^nT_v\) 中落在 \(\mathcal X_E\) 上的相应剩余类一一对应。故可先定义
\[
\widetilde F_0:\pi_n(\mathcal X_E)\to A,
\]
再把 \(\widetilde F_0\) 向整个有限商 \(T_v/B^nT_v\) 任意延拓为 \(\widetilde F\)，于是满足
\[
F=\widetilde F\circ \pi_n|_{\mathcal X_E}.
\]
证毕。
\end{proof}

\begin{theorem}[有限 Grothendieck 有理秩目标中不存在乘法对象的同态嵌入]\label{thm:xi-time-part62ca-no-multiplicative-embedding-into-finite-grothendieck-rank}
\leanverified{paper\_xi\_time\_part62ca\_no\_multiplicative\_embedding\_into\_finite\_grothendieck\_rank}
设 \(M\) 为交换可消去幺半群，并满足
\[
\operatorname{rank}_{\QQ}\!\bigl(G(M)\otimes_{\ZZ}\QQ\bigr)<\infty.
\]
则不存在幺半群单射
\[
\NN_{\times}\hookrightarrow M.
\]
\end{theorem}
\begin{proof}
反设存在单射幺半群同态
\[
f:\NN_{\times}\hookrightarrow M.
\]
由命题 \ref{prop:killo-grothendieck-completion-preserves-injection}，Grothendieck 完成后得到群单射
\[
G(f):G(\NN_{\times})\hookrightarrow G(M).
\]
而由定理 \ref{thm:killo-prime-freedom-non-finitizability} 的第 \(1\) 条，
\[
G(\NN_{\times})\cong \bigoplus_{p\in\mathbb P}\ZZ.
\]
再与 \(\QQ\) 张量。由于 \(\QQ\) 作为 \(\ZZ\)-模是平坦的，单射 \(G(f)\) 仍保持为单射，从而得到
\[
G(\NN_{\times})\otimes_{\ZZ}\QQ
\cong
\bigoplus_{p\in\mathbb P}\QQ,
\]
其 \(\QQ\)-维数可数无穷。于是 \(G(f)\otimes_{\ZZ}\QQ\) 给出一个从可数无穷维 \(\QQ\)-向量空间到
\[
G(M)\otimes_{\ZZ}\QQ
\]
的单射，这与后者具有有限 \(\QQ\)-维数矛盾。故所述单射不存在。证毕。
\end{proof}

\begin{theorem}[Lee--Yang 三次分支不变量共享同一唯一二次分辨子域]\label{thm:xi-time-part62ca-leyang-branch-invariants-common-quadratic-resolvent}
沿用定理 \ref{thm:xi-terminal-zm-stokes-leyang-shared-artin-representation}、推论 \ref{cor:xi-terminal-zm-stokes-leyang-common-quadratic-resolvent} 与定理 \ref{thm:xi-time-part62d-leyang-signed-amplitude-exact-degree6} 的记号，记
\[
u:=\kappa_*^2,
\qquad
b:=B^2,
\qquad
K:=\QQ(u)=\QQ(b)=\QQ(\lambda_*).
\]
设 \(L_K\) 为三次数域 \(K/\QQ\) 的伽罗瓦闭包。则 \(u\)、\(b\) 与 \(\lambda_*\) 的三次最小多项式的分裂域全都等于 \(L_K\)，并且共有同一个唯一非平凡二次子域
\[
\QQ(\sqrt{-111}).
\]
等价地，Lee--Yang 三次核 \(P_{\mathrm{LY}}(y)=256y^3+411y^2+165y+32\) 所定义分裂域的唯一二次分辨子域，恰与三元分支不变量 \((u,b,\lambda_*)\) 的共同二次分辨子域一致。
\end{theorem}
\begin{proof}
由定理 \ref{thm:xi-terminal-zm-stokes-leyang-shared-artin-representation}，
\[
\QQ(u)=\QQ(b)=\QQ(\lambda_*)=K.
\]
因此 \(u\)、\(b\) 与 \(\lambda_*\) 的三次最小多项式都生成同一非伽罗瓦三次数域 \(K\)，其分裂域必同为 \(L_K\)。

再由推论 \ref{cor:xi-terminal-zm-stokes-leyang-common-quadratic-resolvent}，\(L_K/\QQ\) 的唯一二次子域为
\[
\QQ(\sqrt{-111}).
\]
另一方面，定理 \ref{thm:xi-time-part9g-leyang-cubic-discriminant} 已证明 Lee--Yang 三次核 \(P_{\mathrm{LY}}(y)\) 的分裂域同样具有唯一二次子域
\[
\QQ(\sqrt{-111}).
\]
故 Lee--Yang 三次核与分支不变量 \((u,b,\lambda_*)\) 在二次分辨层上严格汇合到同一个域
\[
\QQ(\sqrt{-111}).
\]
证毕。
\end{proof}
